% GENERATED FILE. Edit canonical lessons, then run npm run generate:books.
% canonical-source-hash: fnv1a64:59031bac66cccf98
% canonical-lessons: KA-C444-anusarisu, KA-C444-ullanghisu, KA-C444-palisu, KA-C444-nisedhisu, KA-C444-ekamukha

\chapter{Rules of the Road}
\label{ch:ka-rules-of-the-road}
\hlchaptermodality{444}

\begin{chapteropening}
\textbf{By the end of this chapter:} \emph{I can say to comply with, to violate, to obey, to ban and one-way.}

Complete the last lesson of chapter 444: I can say to comply with, to violate, to obey, to ban and one-way.
\end{chapteropening}

\section[\texorpdfstring{\kn{ಅನುಸರಿಸು} (anusarisu)}{ಅನುಸರಿಸು (anusarisu)}]{\kn{ಅನುಸರಿಸು} (anusarisu) — to comply with; to go by}

\label{lesson:KA-C444-anusarisu}



\begin{quote}
Before the new one: say the Kannada for congestion, then the Kannada for to control.
\end{quote}

\subsection*{You'll want to know: \kn{ಅನುಸರಿಸು}}
\textbf{\kn{ಅನುಸರಿಸು}} — \emph{anusarisu} — \textquotedblleft{}to comply with; to go by\textquotedblright{}.

\begin{quote}
\kn{ಎಲ್ಲಾ} \kn{ಚಾಲಕರು} \kn{ಸಿಗ್ನಲ್} \kn{ಅನುಸರಿಸಬೇಕು}.
\end{quote}

\begin{quote}
\emph{ellā cālakaru signal anusarisabēku.}
\end{quote}

\begin{quote}
All drivers must go by the signals.
\end{quote}

Past (I …): \kn{ಅನುಸರಿಸಿದೆ}, \emph{anusariside}.

\kn{ನಿಯಮ} \kn{ಅನುಸರಿಸು}, to follow a rule; \kn{ಹಿಂಬಾಲಿಸು} is to follow behind someone.

\begin{cousinweb}
\textbf{\kn{ಅನುಸರಿಸು}} (\emph{anusarisu}) is built from Sanskrit \emph{anusaraṇa}, \textquotedblleft{}going after\textquotedblright{}, the noun of \emph{anusarati}, \textquotedblleft{}runs after, follows\textquotedblright{} (CDIAL 338), with the verb ending -\emph{isu}.
\end{cousinweb}

\begin{grammarlens}[title={What you've built}]
One more word for getting around.
\end{grammarlens}

\subsection*{Your turn}
\begin{itemize}
\raggedright
  \item \emph{Say it:} \emph{anusarisu}
  \item \emph{Say it:} \emph{anusarisu}, once more
  \item \emph{Say it:} \emph{ellā cālakaru signal anusarisabēku.} — All drivers must go by the signals.
  \item \emph{Recall:} say the Kannada for congestion, then the Kannada for to control, then say \emph{anusarisu} again
\end{itemize}

\begin{tcolorbox}[breakable,colback=teal!4,colframe=teal!35!black,title={Before you move on}]
What does \kn{ಅನುಸರಿಸು} mean? (\textquotedblleft{}To comply with; to go by\textquotedblright{}.) Say it once more.
\end{tcolorbox}
\section[\texorpdfstring{\kn{ಉಲ್ಲಂಘಿಸು} (ullaṅghisu)}{ಉಲ್ಲಂಘಿಸು (ullaṅghisu)}]{\kn{ಉಲ್ಲಂಘಿಸು} (ullaṅghisu) — to violate; to break a rule}

\label{lesson:KA-C444-ullanghisu}



\begin{quote}
Before the new one: say the Kannada for to control, then the Kannada for to comply with.
\end{quote}

\subsection*{You'll want to know: \kn{ಉಲ್ಲಂಘಿಸು}}
\textbf{\kn{ಉಲ್ಲಂಘಿಸು}} — \emph{ullaṅghisu} — \textquotedblleft{}to violate; to break a rule\textquotedblright{}.

\begin{quote}
\kn{ಅವನು} \kn{ಕೆಂಪು} \kn{ಸಿಗ್ನಲ್} \kn{ಉಲ್ಲಂಘಿಸಿದನು}.
\end{quote}

\begin{quote}
\emph{avanu kempu signal ullaṅghisidanu.}
\end{quote}

\begin{quote}
He jumped the red light.
\end{quote}

Past (I …): \kn{ಉಲ್ಲಂಘಿಸಿದೆ}, \emph{ullaṅghiside}.

\kn{ನಿಯಮ} \kn{ಉಲ್ಲಂಘನೆ}, a breach of the rules.

\begin{cousinweb}
\textbf{\kn{ಉಲ್ಲಂಘಿಸು}} (\emph{ullaṅghisu}) is built from Sanskrit \emph{ullaṅghana}, \textquotedblleft{}leaping over, transgressing\textquotedblright{}, from \emph{ud-} \textquotedblleft{}up, over\textquotedblright{} and \emph{laṅgh}, \textquotedblleft{}to leap over\textquotedblright{} (CDIAL 10905), with the verb ending -\emph{isu}.
\end{cousinweb}

\begin{grammarlens}[title={What you've built}]
One more word for getting around.
\end{grammarlens}

\subsection*{Your turn}
\begin{itemize}
\raggedright
  \item \emph{Say it:} \emph{ullaṅghisu}
  \item \emph{Say it:} \emph{ullaṅghisu}, once more
  \item \emph{Say it:} \emph{avanu kempu signal ullaṅghisidanu.} — He jumped the red light.
  \item \emph{Recall:} say the Kannada for to control, then the Kannada for to comply with, then say \emph{ullaṅghisu} again
\end{itemize}

\begin{tcolorbox}[breakable,colback=teal!4,colframe=teal!35!black,title={Before you move on}]
What does \kn{ಉಲ್ಲಂಘಿಸು} mean? (\textquotedblleft{}To violate; to break a rule\textquotedblright{}.) Say it once more.
\end{tcolorbox}
\section[\texorpdfstring{\kn{ಪಾಲಿಸು} (pālisu)}{ಪಾಲಿಸು (pālisu)}]{\kn{ಪಾಲಿಸು} (pālisu) — to obey; to keep (a rule)}

\label{lesson:KA-C444-palisu}



\begin{quote}
Before the new one: say the Kannada for to comply with, then the Kannada for to violate.
\end{quote}

\subsection*{You'll want to know: \kn{ಪಾಲಿಸು}}
\textbf{\kn{ಪಾಲಿಸು}} — \emph{pālisu} — \textquotedblleft{}to obey; to keep (a rule)\textquotedblright{}.

\begin{quote}
\kn{ನಾವು} \kn{ಸಂಚಾರ} \kn{ನಿಯಮಗಳನ್ನು} \kn{ಪಾಲಿಸಿದೆವು}.
\end{quote}

\begin{quote}
\emph{nāvu sañcāra niyamagaḷannu pālisidevu.}
\end{quote}

\begin{quote}
We kept to the traffic rules.
\end{quote}

Past (I …): \kn{ಪಾಲಿಸಿದೆ}, \emph{pāliside}.

Also to look after or rear; \kn{ನಿಯಮ} \kn{ಪಾಲಿಸು} is the usual pair.

\begin{cousinweb}
\textbf{\kn{ಪಾಲಿಸು}} (\emph{pālisu}) is built from Sanskrit \emph{pālayati}, \textquotedblleft{}guards, protects\textquotedblright{} (CDIAL 8129), with the Kannada verb ending -\emph{isu}: to keep a rule is to guard it.
\end{cousinweb}

\begin{grammarlens}[title={What you've built}]
One more word for getting around.
\end{grammarlens}

\subsection*{Your turn}
\begin{itemize}
\raggedright
  \item \emph{Say it:} \emph{pālisu}
  \item \emph{Say it:} \emph{pālisu}, once more
  \item \emph{Say it:} \emph{nāvu sañcāra niyamagaḷannu pālisidevu.} — We kept to the traffic rules.
  \item \emph{Recall:} say the Kannada for to comply with, then the Kannada for to violate, then say \emph{pālisu} again
\end{itemize}

\begin{tcolorbox}[breakable,colback=teal!4,colframe=teal!35!black,title={Before you move on}]
What does \kn{ಪಾಲಿಸು} mean? (\textquotedblleft{}To obey; to keep (a rule)\textquotedblright{}.) Say it once more.
\end{tcolorbox}
\section[\texorpdfstring{\kn{ನಿಷೇಧಿಸು} (niṣēdhisu)}{ನಿಷೇಧಿಸು (niṣēdhisu)}]{\kn{ನಿಷೇಧಿಸು} (niṣēdhisu) — to ban; to prohibit}

\label{lesson:KA-C444-nisedhisu}



\begin{quote}
Before the new one: say the Kannada for to violate, then the Kannada for to obey.
\end{quote}

\subsection*{You'll want to know: \kn{ನಿಷೇಧಿಸು}}
\textbf{\kn{ನಿಷೇಧಿಸು}} — \emph{niṣēdhisu} — \textquotedblleft{}to ban; to prohibit\textquotedblright{}.

\begin{quote}
\kn{ಉದ್ಯಾನದಲ್ಲಿ} \kn{ಪ್ಲಾಸ್ಟಿಕ್} \kn{ಚೀಲಗಳನ್ನು} \kn{ನಿಷೇಧಿಸಿದ್ದಾರೆ}.
\end{quote}

\begin{quote}
\emph{udyānadalli plāsṭik cīlagaḷannu niṣēdhisiddāre.}
\end{quote}

\begin{quote}
Plastic bags have been banned in the park.
\end{quote}

Past (I …): \kn{ನಿಷೇಧಿಸಿದೆ}, \emph{niṣēdhiside}.

\kn{ನಿಷೇಧ} is a ban: \kn{ಪ್ರವೇಶ} \kn{ನಿಷೇಧ}, no entry.

\begin{cousinweb}
\textbf{\kn{ನಿಷೇಧಿಸು}} (\emph{niṣēdhisu}) is built from Sanskrit \emph{niṣēdha}, \textquotedblleft{}keeping off, prohibition\textquotedblright{} (Monier-Williams), with the verb ending -\emph{isu}.
\end{cousinweb}

\begin{grammarlens}[title={What you've built}]
One more word for getting around.
\end{grammarlens}

\subsection*{Your turn}
\begin{itemize}
\raggedright
  \item \emph{Say it:} \emph{niṣēdhisu}
  \item \emph{Say it:} \emph{niṣēdhisu}, once more
  \item \emph{Say it:} \emph{udyānadalli plāsṭik cīlagaḷannu niṣēdhisiddāre.} — Plastic bags have been banned in the park.
  \item \emph{Recall:} say the Kannada for to violate, then the Kannada for to obey, then say \emph{niṣēdhisu} again
\end{itemize}

\begin{tcolorbox}[breakable,colback=teal!4,colframe=teal!35!black,title={Before you move on}]
What does \kn{ನಿಷೇಧಿಸು} mean? (\textquotedblleft{}To ban; to prohibit\textquotedblright{}.) Say it once more.
\end{tcolorbox}
\section[\texorpdfstring{\kn{ಏಕಮುಖ} (ēkamukha)}{ಏಕಮುಖ (ēkamukha)}]{\kn{ಏಕಮುಖ} (ēkamukha) — one-way}

\label{lesson:KA-C444-ekamukha}



\begin{quote}
Before the new one: say the Kannada for to obey, then the Kannada for to ban.
\end{quote}

\subsection*{You'll want to know: \kn{ಏಕಮುಖ}}
\textbf{\kn{ಏಕಮುಖ}} — \emph{ēkamukha} — \textquotedblleft{}one-way\textquotedblright{}.

\begin{quote}
\kn{ಇದು} \kn{ಏಕಮುಖ} \kn{ರಸ್ತೆ}, \kn{ಆ} \kn{ಕಡೆಯಿಂದ} \kn{ಬರಬಾರದು}.
\end{quote}

\begin{quote}
\emph{idu ēkamukha raste, ā kaḍeyinda barabāradu.}
\end{quote}

\begin{quote}
This is a one-way road; nobody may come from that side.
\end{quote}

\kn{ಏಕ}, one, plus \kn{ಮುಖ}, face: facing one way.

\begin{cousinweb}
\textbf{\kn{ಏಕಮುಖ}} (\emph{ēkamukha}) is built from two Sanskrit words: \emph{ēka}, \textquotedblleft{}one\textquotedblright{} (CDIAL 2462), and \emph{mukha}, \textquotedblleft{}mouth, face, entrance\textquotedblright{} (CDIAL 10158).
\end{cousinweb}

\begin{grammarlens}[title={What you've built}]
One more word for getting around.
\end{grammarlens}

\subsection*{Your turn}
\begin{itemize}
\raggedright
  \item \emph{Say it:} \emph{ēkamukha}
  \item \emph{Say it:} \emph{ēkamukha}, once more
  \item \emph{Say it:} \emph{idu ēkamukha raste, ā kaḍeyinda barabāradu.} — This is a one-way road; nobody may come from that side.
  \item \emph{Recall:} say the Kannada for to obey, then the Kannada for to ban, then say \emph{ēkamukha} again
\end{itemize}

\begin{tcolorbox}[breakable,colback=teal!4,colframe=teal!35!black,title={Before you move on}]
What does \kn{ಏಕಮುಖ} mean? (\textquotedblleft{}One-way\textquotedblright{}.) Say it once more.
\end{tcolorbox}
